% =====================================================================
% Note on: Hedging the Treasury Lock — Mario Pucci, 2019
% ssrn:3386521
% =====================================================================
\documentclass[11pt,a4paper]{article}

\newcommand{\papertag}{pucci\_2019\_tlock}

% =====================================================================
% Shared preamble for paper notes
% Include with:  % =====================================================================
% Shared preamble for paper notes
% Include with:  % =====================================================================
% Shared preamble for paper notes
% Include with:  \input{../_common/preamble}
% =====================================================================

% --- Core packages ---
\usepackage[utf8]{inputenc}
\usepackage[T1]{fontenc}
\usepackage{lmodern}
\usepackage[margin=2.5cm]{geometry}
\usepackage{amsmath,amssymb,amsthm,mathtools}
\usepackage{bm}
\usepackage{booktabs}
\usepackage{array}
\usepackage{enumitem}
\usepackage{hyperref}
\usepackage{xcolor}
\usepackage{listings}
\usepackage{fancyhdr}
\usepackage{titlesec}

% --- Hyperref styling ---
\hypersetup{
  colorlinks=true,
  linkcolor=blue!50!black,
  urlcolor=blue!50!black,
  citecolor=blue!50!black
}

% --- Theorem-like environments ---
\theoremstyle{definition}
\newtheorem{defn}{Definition}
\newtheorem{remark}{Remark}
\newtheorem{example}{Example}

\theoremstyle{plain}
\newtheorem{prop}{Proposition}
\newtheorem{lemma}{Lemma}

% --- Coloured boxes for the structured sections ---
% Validation block, commentary, TODOs use coloured frames so the eye
% can find them when flipping through the note.
\usepackage[most]{tcolorbox}

\newtcolorbox{commentary}{
  colback=yellow!5,
  colframe=yellow!50!black,
  title=Commentary,
  fonttitle=\bfseries,
  breakable
}

\newtcolorbox{validation}{
  colback=green!3,
  colframe=green!50!black,
  title=Validation,
  fonttitle=\bfseries,
  breakable
}

\newtcolorbox{todobox}{
  colback=red!3,
  colframe=red!50!black,
  title=TODO / Open questions,
  fonttitle=\bfseries,
  breakable
}

% --- Code listing style for Python snippets in validation blocks ---
\lstdefinestyle{py}{
  language=Python,
  basicstyle=\ttfamily\footnotesize,
  keywordstyle=\color{blue!60!black},
  commentstyle=\color{gray},
  stringstyle=\color{red!60!black},
  showstringspaces=false,
  breaklines=true,
  frame=single,
  framesep=4pt,
  rulecolor=\color{gray!30}
}

% --- Page header: shows paper tag so a printed note is identifiable ---
\pagestyle{fancy}
\fancyhf{}
\fancyhead[L]{\small\textit{\papertag}}
\fancyhead[R]{\small\thepage}
\renewcommand{\headrulewidth}{0.4pt}

% --- Section spacing: tighter than article default ---
\titlespacing*{\section}{0pt}{1.5ex plus 0.5ex minus 0.2ex}{1ex plus 0.2ex}
\titlespacing*{\subsection}{0pt}{1.2ex plus 0.3ex minus 0.2ex}{0.6ex plus 0.2ex}

% =====================================================================

% --- Core packages ---
\usepackage[utf8]{inputenc}
\usepackage[T1]{fontenc}
\usepackage{lmodern}
\usepackage[margin=2.5cm]{geometry}
\usepackage{amsmath,amssymb,amsthm,mathtools}
\usepackage{bm}
\usepackage{booktabs}
\usepackage{array}
\usepackage{enumitem}
\usepackage{hyperref}
\usepackage{xcolor}
\usepackage{listings}
\usepackage{fancyhdr}
\usepackage{titlesec}

% --- Hyperref styling ---
\hypersetup{
  colorlinks=true,
  linkcolor=blue!50!black,
  urlcolor=blue!50!black,
  citecolor=blue!50!black
}

% --- Theorem-like environments ---
\theoremstyle{definition}
\newtheorem{defn}{Definition}
\newtheorem{remark}{Remark}
\newtheorem{example}{Example}

\theoremstyle{plain}
\newtheorem{prop}{Proposition}
\newtheorem{lemma}{Lemma}

% --- Coloured boxes for the structured sections ---
% Validation block, commentary, TODOs use coloured frames so the eye
% can find them when flipping through the note.
\usepackage[most]{tcolorbox}

\newtcolorbox{commentary}{
  colback=yellow!5,
  colframe=yellow!50!black,
  title=Commentary,
  fonttitle=\bfseries,
  breakable
}

\newtcolorbox{validation}{
  colback=green!3,
  colframe=green!50!black,
  title=Validation,
  fonttitle=\bfseries,
  breakable
}

\newtcolorbox{todobox}{
  colback=red!3,
  colframe=red!50!black,
  title=TODO / Open questions,
  fonttitle=\bfseries,
  breakable
}

% --- Code listing style for Python snippets in validation blocks ---
\lstdefinestyle{py}{
  language=Python,
  basicstyle=\ttfamily\footnotesize,
  keywordstyle=\color{blue!60!black},
  commentstyle=\color{gray},
  stringstyle=\color{red!60!black},
  showstringspaces=false,
  breaklines=true,
  frame=single,
  framesep=4pt,
  rulecolor=\color{gray!30}
}

% --- Page header: shows paper tag so a printed note is identifiable ---
\pagestyle{fancy}
\fancyhf{}
\fancyhead[L]{\small\textit{\papertag}}
\fancyhead[R]{\small\thepage}
\renewcommand{\headrulewidth}{0.4pt}

% --- Section spacing: tighter than article default ---
\titlespacing*{\section}{0pt}{1.5ex plus 0.5ex minus 0.2ex}{1ex plus 0.2ex}
\titlespacing*{\subsection}{0pt}{1.2ex plus 0.3ex minus 0.2ex}{0.6ex plus 0.2ex}

% =====================================================================

% --- Core packages ---
\usepackage[utf8]{inputenc}
\usepackage[T1]{fontenc}
\usepackage{lmodern}
\usepackage[margin=2.5cm]{geometry}
\usepackage{amsmath,amssymb,amsthm,mathtools}
\usepackage{bm}
\usepackage{booktabs}
\usepackage{array}
\usepackage{enumitem}
\usepackage{hyperref}
\usepackage{xcolor}
\usepackage{listings}
\usepackage{fancyhdr}
\usepackage{titlesec}

% --- Hyperref styling ---
\hypersetup{
  colorlinks=true,
  linkcolor=blue!50!black,
  urlcolor=blue!50!black,
  citecolor=blue!50!black
}

% --- Theorem-like environments ---
\theoremstyle{definition}
\newtheorem{defn}{Definition}
\newtheorem{remark}{Remark}
\newtheorem{example}{Example}

\theoremstyle{plain}
\newtheorem{prop}{Proposition}
\newtheorem{lemma}{Lemma}

% --- Coloured boxes for the structured sections ---
% Validation block, commentary, TODOs use coloured frames so the eye
% can find them when flipping through the note.
\usepackage[most]{tcolorbox}

\newtcolorbox{commentary}{
  colback=yellow!5,
  colframe=yellow!50!black,
  title=Commentary,
  fonttitle=\bfseries,
  breakable
}

\newtcolorbox{validation}{
  colback=green!3,
  colframe=green!50!black,
  title=Validation,
  fonttitle=\bfseries,
  breakable
}

\newtcolorbox{todobox}{
  colback=red!3,
  colframe=red!50!black,
  title=TODO / Open questions,
  fonttitle=\bfseries,
  breakable
}

% --- Code listing style for Python snippets in validation blocks ---
\lstdefinestyle{py}{
  language=Python,
  basicstyle=\ttfamily\footnotesize,
  keywordstyle=\color{blue!60!black},
  commentstyle=\color{gray},
  stringstyle=\color{red!60!black},
  showstringspaces=false,
  breaklines=true,
  frame=single,
  framesep=4pt,
  rulecolor=\color{gray!30}
}

% --- Page header: shows paper tag so a printed note is identifiable ---
\pagestyle{fancy}
\fancyhf{}
\fancyhead[L]{\small\textit{\papertag}}
\fancyhead[R]{\small\thepage}
\renewcommand{\headrulewidth}{0.4pt}

% --- Section spacing: tighter than article default ---
\titlespacing*{\section}{0pt}{1.5ex plus 0.5ex minus 0.2ex}{1ex plus 0.2ex}
\titlespacing*{\subsection}{0pt}{1.2ex plus 0.3ex minus 0.2ex}{0.6ex plus 0.2ex}

% =====================================================================
% House notation: shared symbols used across all paper notes.
% Each note imports this and may shadow individual macros if a paper
% uses a clashing convention — but the *default* meaning is fixed here.
%
% Include with:  % =====================================================================
% House notation: shared symbols used across all paper notes.
% Each note imports this and may shadow individual macros if a paper
% uses a clashing convention — but the *default* meaning is fixed here.
%
% Include with:  % =====================================================================
% House notation: shared symbols used across all paper notes.
% Each note imports this and may shadow individual macros if a paper
% uses a clashing convention — but the *default* meaning is fixed here.
%
% Include with:  \input{../_common/notation}
% =====================================================================

% --- Measures and expectations ---
\newcommand{\Q}{\mathbb{Q}}              % risk-neutral measure
\newcommand{\Qt}{\mathbb{Q}^{T}}         % T-forward measure
\newcommand{\E}{\mathbb{E}}              % expectation
\newcommand{\Et}[1]{\E^{#1}}             % expectation under measure #1
\newcommand{\Ftime}[1]{\mathcal{F}_{#1}} % filtration at #1

% --- Discount and forward objects ---
\newcommand{\Disc}[2]{D_{#1,#2}}         % discount factor D_{t,T}
\newcommand{\Bond}[2]{P_{#1}(#2)}        % bond price P_t(y)
\newcommand{\Ann}[1]{A_{#1}}             % annuity / level
\newcommand{\Numer}[1]{N_{#1}}           % numeraire
\newcommand{\Fwd}[2]{F_{#1,#2}}          % forward price F_{t,T}
\newcommand{\Fwdpx}[2]{\mathrm{Fwd}_{#1}(#2)} % verbose forward-price F_t(T)

% --- Rates ---
\newcommand{\IRR}{\mathrm{IRR}}          % internal rate of return
\newcommand{\Yld}{y}                     % bond yield variable
\newcommand{\Strike}{K}                  % strike (price-space)
\newcommand{\Lock}{L}                    % locked yield / rate
\newcommand{\Repo}{r^{\mathrm{repo}}}    % repo rate
\newcommand{\Fund}{r^{\mathrm{fun}}}     % unsecured funding rate
\newcommand{\OIS}{r^{\mathrm{ois}}}      % OIS rate
\newcommand{\Coll}{r^{\mathrm{c}}}       % collateral rate
\newcommand{\Rcmt}{R^{\mathrm{cmt}}}     % constant-maturity-treasury rate
\newcommand{\Rswp}{R^{\mathrm{swp}}}     % swap rate (used for risky variant)
\newcommand{\Rec}{\mathrm{Rec}}          % recovery rate
\newcommand{\NPV}{\mathrm{NPV}}          % present-value functional
\newcommand{\PVone}{\mathrm{PV01}}       % risky annuity
\newcommand{\Loss}{\mathrm{Loss}}        % default-leg integral
\newcommand{\Sasw}{S^{\mathrm{ASW}}}     % asset-swap spread
\newcommand{\Scds}{S^{\mathrm{CDS}}}     % CDS spread
\newcommand{\Basis}{\mathrm{Basis}}      % CDS-bond basis

% --- Sensitivities ---
\newcommand{\Diff}[2]{D_{#1}^{#2}}       % Euler's notation: D_x^k[f]
\newcommand{\Riskf}{\mathrm{RiskFactor}} % bond risk factor (= -DV01*1e4)
\newcommand{\DVone}{\mathrm{DV01}}

% --- Indicator / misc ---
\newcommand{\Ind}{\mathbf{1}}
\newcommand{\given}{\,|\,}
\newcommand{\dd}{\,\mathrm{d}}

% --- Greeks-style ---
\newcommand{\Gam}{\Gamma}
\newcommand{\Vega}{\mathcal{V}}

% =====================================================================

% --- Measures and expectations ---
\newcommand{\Q}{\mathbb{Q}}              % risk-neutral measure
\newcommand{\Qt}{\mathbb{Q}^{T}}         % T-forward measure
\newcommand{\E}{\mathbb{E}}              % expectation
\newcommand{\Et}[1]{\E^{#1}}             % expectation under measure #1
\newcommand{\Ftime}[1]{\mathcal{F}_{#1}} % filtration at #1

% --- Discount and forward objects ---
\newcommand{\Disc}[2]{D_{#1,#2}}         % discount factor D_{t,T}
\newcommand{\Bond}[2]{P_{#1}(#2)}        % bond price P_t(y)
\newcommand{\Ann}[1]{A_{#1}}             % annuity / level
\newcommand{\Numer}[1]{N_{#1}}           % numeraire
\newcommand{\Fwd}[2]{F_{#1,#2}}          % forward price F_{t,T}
\newcommand{\Fwdpx}[2]{\mathrm{Fwd}_{#1}(#2)} % verbose forward-price F_t(T)

% --- Rates ---
\newcommand{\IRR}{\mathrm{IRR}}          % internal rate of return
\newcommand{\Yld}{y}                     % bond yield variable
\newcommand{\Strike}{K}                  % strike (price-space)
\newcommand{\Lock}{L}                    % locked yield / rate
\newcommand{\Repo}{r^{\mathrm{repo}}}    % repo rate
\newcommand{\Fund}{r^{\mathrm{fun}}}     % unsecured funding rate
\newcommand{\OIS}{r^{\mathrm{ois}}}      % OIS rate
\newcommand{\Coll}{r^{\mathrm{c}}}       % collateral rate
\newcommand{\Rcmt}{R^{\mathrm{cmt}}}     % constant-maturity-treasury rate
\newcommand{\Rswp}{R^{\mathrm{swp}}}     % swap rate (used for risky variant)
\newcommand{\Rec}{\mathrm{Rec}}          % recovery rate
\newcommand{\NPV}{\mathrm{NPV}}          % present-value functional
\newcommand{\PVone}{\mathrm{PV01}}       % risky annuity
\newcommand{\Loss}{\mathrm{Loss}}        % default-leg integral
\newcommand{\Sasw}{S^{\mathrm{ASW}}}     % asset-swap spread
\newcommand{\Scds}{S^{\mathrm{CDS}}}     % CDS spread
\newcommand{\Basis}{\mathrm{Basis}}      % CDS-bond basis

% --- Sensitivities ---
\newcommand{\Diff}[2]{D_{#1}^{#2}}       % Euler's notation: D_x^k[f]
\newcommand{\Riskf}{\mathrm{RiskFactor}} % bond risk factor (= -DV01*1e4)
\newcommand{\DVone}{\mathrm{DV01}}

% --- Indicator / misc ---
\newcommand{\Ind}{\mathbf{1}}
\newcommand{\given}{\,|\,}
\newcommand{\dd}{\,\mathrm{d}}

% --- Greeks-style ---
\newcommand{\Gam}{\Gamma}
\newcommand{\Vega}{\mathcal{V}}

% =====================================================================

% --- Measures and expectations ---
\newcommand{\Q}{\mathbb{Q}}              % risk-neutral measure
\newcommand{\Qt}{\mathbb{Q}^{T}}         % T-forward measure
\newcommand{\E}{\mathbb{E}}              % expectation
\newcommand{\Et}[1]{\E^{#1}}             % expectation under measure #1
\newcommand{\Ftime}[1]{\mathcal{F}_{#1}} % filtration at #1

% --- Discount and forward objects ---
\newcommand{\Disc}[2]{D_{#1,#2}}         % discount factor D_{t,T}
\newcommand{\Bond}[2]{P_{#1}(#2)}        % bond price P_t(y)
\newcommand{\Ann}[1]{A_{#1}}             % annuity / level
\newcommand{\Numer}[1]{N_{#1}}           % numeraire
\newcommand{\Fwd}[2]{F_{#1,#2}}          % forward price F_{t,T}
\newcommand{\Fwdpx}[2]{\mathrm{Fwd}_{#1}(#2)} % verbose forward-price F_t(T)

% --- Rates ---
\newcommand{\IRR}{\mathrm{IRR}}          % internal rate of return
\newcommand{\Yld}{y}                     % bond yield variable
\newcommand{\Strike}{K}                  % strike (price-space)
\newcommand{\Lock}{L}                    % locked yield / rate
\newcommand{\Repo}{r^{\mathrm{repo}}}    % repo rate
\newcommand{\Fund}{r^{\mathrm{fun}}}     % unsecured funding rate
\newcommand{\OIS}{r^{\mathrm{ois}}}      % OIS rate
\newcommand{\Coll}{r^{\mathrm{c}}}       % collateral rate
\newcommand{\Rcmt}{R^{\mathrm{cmt}}}     % constant-maturity-treasury rate
\newcommand{\Rswp}{R^{\mathrm{swp}}}     % swap rate (used for risky variant)
\newcommand{\Rec}{\mathrm{Rec}}          % recovery rate
\newcommand{\NPV}{\mathrm{NPV}}          % present-value functional
\newcommand{\PVone}{\mathrm{PV01}}       % risky annuity
\newcommand{\Loss}{\mathrm{Loss}}        % default-leg integral
\newcommand{\Sasw}{S^{\mathrm{ASW}}}     % asset-swap spread
\newcommand{\Scds}{S^{\mathrm{CDS}}}     % CDS spread
\newcommand{\Basis}{\mathrm{Basis}}      % CDS-bond basis

% --- Sensitivities ---
\newcommand{\Diff}[2]{D_{#1}^{#2}}       % Euler's notation: D_x^k[f]
\newcommand{\Riskf}{\mathrm{RiskFactor}} % bond risk factor (= -DV01*1e4)
\newcommand{\DVone}{\mathrm{DV01}}

% --- Indicator / misc ---
\newcommand{\Ind}{\mathbf{1}}
\newcommand{\given}{\,|\,}
\newcommand{\dd}{\,\mathrm{d}}

% --- Greeks-style ---
\newcommand{\Gam}{\Gamma}
\newcommand{\Vega}{\mathcal{V}}


\title{Note on: \emph{Hedging the Treasury Lock}\\
       \large Mario Pucci (Banca IMI), 2019}
\author{Reading notes}
\date{\today}

\begin{document}
\maketitle

% =====================================================================
\section*{Metadata}
\begin{tabular}{@{}p{3.5cm}p{11cm}@{}}
\toprule
Title           & Hedging the Treasury Lock \\
Author(s)       & Mario Pucci (Banca IMI SpA) \\
Year / version  & 19 August 2019 \\
Source          & \href{https://ssrn.com/abstract=3386521}{ssrn:3386521} \\
Local file      & \texttt{papers/pucci\_2019\_tlock/ssrn3386521.pdf} \\
Tags            & \texttt{treasury}, \texttt{forward}, \texttt{repo},
                  \texttt{convexity}, \texttt{pre-hedge}, \texttt{hedge-accounting} \\
Date read       & \today \\
\bottomrule
\end{tabular}

% =====================================================================
\section*{One-paragraph summary}
A T-Lock pays the client $N \cdot \Riskf_{t_e} \cdot (\IRR_{t_e} - \Lock)$
at expiry $t_e$, where $\IRR$ is the bond's internal rate of return and
$\Riskf = -\partial P/\partial y \cdot 10^4$ (i.e.\ $-10^4 \cdot \DVone$).
Because $P(y)$ is convex, the linear payoff in $\IRR$ is, in fact,
\emph{globally below} the negative of $P$ when viewed as a function of
yield --- so booking a short T-Lock as a long forward contract on the
same bond, struck at $\Bond{t_e}{\Lock}$, is an overhedge. The paper
quantifies that overhedge error, shows the gamma is locally negative
(so delta-hedging the proxy forward generates gamma profits), and
extends the argument to Then-Current-T-Locks via a benchmark-roll P\&L
analysis backed by 2008--2019 Treasury auction statistics. The natural
hedge is a reverse repo on the underlying bond; the bond-specific repo
rate is the right drift in the forward price.

% =====================================================================
\section{Setup and notation}

The paper's symbols mostly align with the house notation, with one
collision worth flagging: $P_t(y)$ in the paper is the \emph{bond}
price as a function of yield, not a discount factor. House notation
keeps $\Bond{t}{y}$ for the bond and $\Disc{t}{T}$ for the discount
factor; they should not be confused.

\begin{tabular}{@{}p{3cm}p{3.5cm}p{8.5cm}@{}}
\toprule
Paper & House & Meaning \\
\midrule
$P_t(y)$       & $\Bond{t}{y}$        & Dirty price of the bond at time $t$ as a function of yield $y$ \\
$P_t^{mkt}$    & $\Bond{t}{\IRR_t}$   & Observed market dirty price \\
$\IRR_t$       & $\IRR_t$             & Internal rate of return solving $\Bond{t}{\IRR_t} = P_t^{mkt}$ \\
$L$            & $\Lock$              & Locked yield (strike) \\
$\Riskf_t$     & $\Riskf_t$           & $-\partial \Bond{t}{y}/\partial y$ at $y=\IRR_t$. Equals $-10^4 \cdot \DVone$. \\
$D_{t\,t_e}$   & $\Disc{t}{t_e}$      & Risk-free DF from $t$ to $t_e$ \\
$\mathrm{ForwardPrice}_t(t_e)$ & $\Fwdpx{t}{t_e}$ & Forward bond price for delivery at $t_e$ \\
$r^{\mathrm{repo}}_{tt_e}$ & $\Repo_{t,t_e}$ & Simply-compounded forward repo rate $t \to t_e$ \\
$a$            & $a \in \{+1,-1\}$    & Long ($+1$) / short ($-1$) T-Lock \\
$P_t^{(0)}$    & ---                  & Bond price stripped of coupons ($c=0$); zero-bond proxy \\
$A_t(y)$       & $\Ann{t}(y)$         & Coupon annuity $\sum \alpha_i e^{-(t_i-t)y}$ \\
\bottomrule
\end{tabular}

% =====================================================================
\section{Key equations}

\paragraph{Payoff and IRR definition.}
\begin{equation}\label{eq:payoff}
  g \;=\; a \cdot \Riskf_{t_e} \cdot (\IRR_{t_e} - \Lock).
\end{equation}
\textit{Says:} a long ($a=1$) T-Lock pays the bondholder when yields
rise (because $\Riskf>0$); a short pays when yields fall.

\begin{equation}\label{eq:irr-def}
  \Bond{t}{\IRR_t} \;=\; P_t^{mkt}, \qquad
  \Bond{t}{y} \;=\; \prod_{i=1}^n \frac{1}{1+\alpha_i y}
   + c \sum_{i=1}^n \alpha_i \prod_{j=1}^i \frac{1}{1+\alpha_j y}.
\end{equation}
\textit{Says:} the IRR is the yield that reprices the bond to its
dirty market price; $\Bond{t}{\cdot}$ uses the \emph{compounded} yield
convention (not continuously compounded --- this matters for code).

\paragraph{Continuous-compounding form (used for greeks).}
For the sensitivity analysis, the paper switches to
\begin{equation}\label{eq:bond-cc}
  \Bond{t}{y} \;=\; e^{-(t_n - t)y} \;+\; c \sum_{t_i > t} \alpha_i e^{-(t_i-t)y},
\end{equation}
with derivatives (B.84--87)
\begin{equation}\label{eq:bond-derivs}
  \Diff{y}{k}[\Bond{t}{\cdot}] \;=\; (t - t_n)^k e^{-(t_n - t)y}
                             \;+\; c \sum_{t_i > t} \alpha_i (t-t_i)^k e^{-(t_i-t)y}.
\end{equation}
\textit{Says:} closed-form derivatives of all orders in $y$, sign
alternating with $k$. $\Diff{y}{1}<0$, $\Diff{y}{2}>0$, $\Diff{y}{3}<0$.

\paragraph{Linear (forward) approximation.}
\begin{equation}\label{eq:linear-approx}
  g \;\simeq\; a \cdot \bigl(\Bond{t_e}{\Lock} - \Bond{t_e}{\IRR_{t_e}}\bigr)
       \;=\; a \cdot \bigl(\Bond{t_e}{\Lock} - P_{t_e}^{mkt}\bigr).
\end{equation}
\textit{Says:} at first order, the T-Lock is a forward contract on the
bond struck at $\Strike = \Bond{t_e}{\Lock}$. Booking a short T-Lock
as a long forward at this strike is the operational shortcut.

\paragraph{Overhedge error bound.}
With the remainder $R_1(y) = \Bond{}{y} - \Bond{}{\Lock} - \Diff{y}{1}[\Bond{}{\cdot}](y-\Lock)$,
\begin{equation}\label{eq:overhedge}
  \lvert R_1(y) \rvert \;\le\; \frac{M\,(y-\Lock)^2}{2},
  \qquad M = \max_I \bigl|\Diff{y}{2}[\Bond{}{\cdot}]\bigr|.
\end{equation}
\textit{Says:} for the 10Y benchmark with $M \approx 150$ and a 30\%
yield change (i.e.\ $|y-\Lock| \approx 0.3\Lock$), the price error is
about \textbf{0.005\%} on 100\% notional. The approximation is tight.

\paragraph{Delta and gamma.}
\begin{equation}\label{eq:delta}
  \Delta \;=\; \Diff{P}{1}[g] \cdot \Disc{t}{t_e}, \qquad
  \Diff{y}{1}[g] \;=\; -a \cdot \bigl[\Diff{y}{2}[\Bond{}{\cdot}](y-\Lock) + \Diff{y}{1}[\Bond{}{\cdot}]\bigr],
\end{equation}
\begin{equation}\label{eq:gamma-local}
  \Gam \;=\; \Diff{P}{2}[g] \;\simeq\; -a \cdot \Diff{y}{2}[\Bond{}{\cdot}] \cdot \bigl[1/\Diff{y}{1}[\Bond{}{\cdot}]\bigr]^2
  \qquad \text{(local, } |y-\Lock| \text{ small)}.
\end{equation}
\textit{Says:} for $a=1$ near $\Lock$, $\Delta < 0$ (consistent with
``long T-Lock = synthetic short bond'') and $\Gam < 0$. The negative
gamma is what makes the linear proxy an overhedge --- the bank gains
on dynamic delta-rebalancing.

\paragraph{Sharp sign conditions} (resolve \eqref{eq:delta} and the gamma directly):
\begin{align}
  \Delta < 0 &\iff y - \Lock \;<\; -\Diff{y}{1}[\Bond{}{\cdot}] \;/\; \Diff{y}{2}[\Bond{}{\cdot}] \;>\; 0, \label{eq:delta-cond}\\
  \Gam   < 0 &\iff y - \Lock \;<\; \frac{\Diff{y}{1}[\Bond{}{\cdot}] \cdot \Diff{y}{3}[\Bond{}{\cdot}]}{(\Diff{y}{2}[\Bond{}{\cdot}])^2 - \Diff{y}{1}[\Bond{}{\cdot}] \cdot \Diff{y}{3}[\Bond{}{\cdot}]} > 0. \label{eq:gamma-cond}
\end{align}
\textit{Says:} both bounds give exact yield thresholds beyond which
the sign flips. In ``normal'' yield regimes ($y$ not far above
$\Lock$), $\Delta < 0$ and $\Gam < 0$ both hold.

\paragraph{Forward price via repo (Appendix D).}
For zero haircut,
\begin{equation}\label{eq:fwd-repo}
  \Fwdpx{t}{t_e} \;=\; P_t^{mkt} \, e^{\Repo \cdot (t_e - t)}
                  \;-\; \sum_{t < t_i \le t_e} c_i \, e^{\Repo \cdot (t_e - t_i)}.
\end{equation}
\textit{Says:} the forward bond price drifts at the \emph{repo} rate,
not the OIS / risk-free rate. Coupons paid before $t_e$ are forwarded
to $t_e$ at the same repo rate. With a haircut $h$ blend in the
unsecured funding rate $\Fund$.

\paragraph{Then-Current-T-Lock: strike at roll.}
\begin{equation}\label{eq:tcurr-strike}
  \Strike \;=\; 1 + \Riskf_{t_e} \cdot (R_t - \Lock),
\end{equation}
\begin{equation}\label{eq:tcurr-strikegap}
  \mathrm{StrikeGap}
   \;=\; (\widehat{\Riskf}_{t_e} - \Riskf_{t_e})\cdot(\widehat{R}_{t_{\mathrm{roll}}} - \Lock)
   \;-\; \Riskf_{t_e}\cdot(R_{t_{\mathrm{roll}}} - \widehat{R}_{t_{\mathrm{roll}}}).
\end{equation}
\textit{Says:} when the on-the-run rolls, the strike updates. The
first term captures the risk-factor jump between old and new
benchmark; the second the IRR jump at the roll date.

% =====================================================================
\section{Derivations \& commentary}

\begin{commentary}
The key trick is \eqref{eq:linear-approx}: turning the T-Lock payoff
$g = -\Delta\Bond{}{\cdot} \cdot (y - \Lock)$ into the secant
$\Bond{t_e}{\Lock} - \Bond{t_e}{\IRR_{t_e}}$ uses the first-order
Taylor identity $\Bond{}{\IRR} - \Bond{}{\Lock} \simeq \Diff{y}{1}[\Bond{}{\cdot}](\IRR - \Lock)$.
The sign of the remainder is fixed by convexity ($\Diff{y}{2}[\Bond{}{\cdot}] > 0$
always for a bond with positive cashflows), which is why the bank's
position is conservative \emph{globally}, not just locally. Importantly,
the gamma being negative is the \emph{opposite} sign to what a naive
``T-Lock $\approx$ CMS-let'' intuition would give --- the CMT
interpretation freezes $\Riskf_{t_e}$ and obscures the dependence of
the risk factor on yield (which is what flips the convexity).
\end{commentary}

\begin{commentary}
The repo rate in \eqref{eq:fwd-repo} is bond-specific. A benchmark
Treasury can trade ``special'' --- the paper notes historical episodes
of the 10Y repo rate going to $-300$~bps. For validation, the test
case in the paper uses GC rates (Tab.~\ref{tab:repo-jan2019}); the
specialness premium must be added as a separate spread input.
\end{commentary}

\begin{commentary}
The continuous-compounding form \eqref{eq:bond-cc} is used purely for
the analytic greeks. The \emph{contract definition} \eqref{eq:irr-def}
uses compounded (i.e.\ $1/(1+\alpha y)$) discounting --- mixing them
in a pricing function is a common bug. House rule: keep two separate
bond-price functions, one per convention, and only the compounded
one defines $\IRR$.
\end{commentary}

% =====================================================================
\section{Validation}

\begin{validation}
This is a \emph{specification} of mathematical objects the paper
defines and the numerical values they should take. It says nothing
about how those objects are implemented or invoked.

\textbf{Quantities the paper defines:}
\begin{itemize}[noitemsep]
  \item $\Bond{t}{y}$, the bond dirty price as a function of yield, in
        two distinct conventions:
        \emph{compounded} (eq.~\ref{eq:irr-def}) and
        \emph{continuously compounded} (eq.~\ref{eq:bond-cc}).
        Only the compounded form defines the contractual $\IRR$.
  \item $\IRR_t$, the unique solution of $\Bond{t}{\IRR_t} = P^{mkt}_t$
        with $\Bond{t}{\cdot}$ in the compounded convention.
  \item $\Diff{y}{k}[\Bond{t}{\cdot}]$, the $k$-th derivative of the
        bond price with respect to yield, in closed form
        (eq.~\ref{eq:bond-derivs}).
  \item $\Riskf_t = -\Diff{y}{1}[\Bond{t}{y}]\big|_{y=\IRR_t}$.
  \item $\Fwdpx{t}{t_e}$, the $t_e$-forward bond price under a repo
        rate (eq.~\ref{eq:fwd-repo}), zero-haircut limit.
  \item $g(y)$, the T-Lock payoff as a function of yield
        (eq.~\ref{eq:payoff}), and its derivatives
        $\Diff{y}{1}[g]$, $\Diff{y}{2}[g]$ in closed form
        (eqs.~\ref{eq:delta}, \ref{eq:gamma-local}).
\end{itemize}

\textbf{Canonical numerical case} (paper, p.~4):
\begin{itemize}[noitemsep]
  \item Instrument: US Treasury, coupon 3.125\%, maturity 25-Nov-2028, ISIN US9128285M81.
  \item As-of: 24 January 2019; $t_e - t = 3$ months.
  \item $P^{mkt} = 104.1055$ (dirty), clean price $= 103.4922$.
  \item $\IRR = 2.717\%$ spot.
  \item 3M GC repo $\Repo = 2.46\%$ (mid of bid/ask, Tab.~\ref{tab:repo-jan2019}).
  \item Forward: $\IRR_{\mathrm{fwd}} = 2.53\%$, dirty $104.74\%$, clean $103.36\%$.
  \item $\Lock = 2.717\%$.
\end{itemize}

\textbf{Expected numerical values:}

\smallskip
\begin{tabular}{@{}p{6cm}rl@{}}
\toprule
Quantity                                                              & Value                      & Tolerance \\
\midrule
$\Fwdpx{t}{t_e}$ (dirty)                                              & $104.74$                   & $\pm 0.01$ \\
$\Fwdpx{t}{t_e}$ (clean)                                              & $103.36$                   & $\pm 0.01$ \\
Forward $\IRR$                                                        & $2.53\%$                   & $\pm 1$~bp \\
$|R_1(y)|$ at 10Y, $|y-\Lock| \approx 0.3\Lock$                       & $\sim 0.005\%$ notional    & order-of-mag. \\
$M = \max_I |\Diff{y}{2}[\Bond{}{\cdot}]|$ (10Y benchmark)            & $\sim 150$                 & order-of-mag. \\
\bottomrule
\end{tabular}
\smallskip

\begin{center}
\textit{Repo rates for the 10Y as of 24-Jan-2019 (Paper Tab.~1)}\\[2pt]
\begin{tabular}{@{}lll@{}}
\toprule
Maturity & bid    & ask    \\
\midrule
1M       & 2.57\% & 2.52\% \\
3M       & 2.60\% & 2.55\% \\
6M       & 2.63\% & 2.58\% \\
12M      & 2.70\% & 2.65\% \\
\bottomrule
\end{tabular}
\refstepcounter{table}\label{tab:repo-jan2019}
\end{center}

\textbf{Equation $\leftrightarrow$ quantity map:}

\smallskip
\begin{tabular}{@{}p{2.8cm}p{3.8cm}p{6.2cm}@{}}
\toprule
Paper eq.                                     & Symbol                                             & Quantity \\
\midrule
\eqref{eq:irr-def}                            & $\Bond{t}{y}$, $\IRR_t$                            & dirty price (compounded); the $y$ for which $\Bond{t}{y}=P^{mkt}_t$ \\
\eqref{eq:bond-cc}                            & $\Bond{t}{y}$ (CC)                                 & dirty price (continuously compounded) \\
\eqref{eq:bond-derivs}                        & $\Diff{y}{k}[\Bond{}{\cdot}]$                      & closed-form $k$-th yield derivative \\
\eqref{eq:linear-approx}                      & $g \simeq a\,(\Bond{t_e}{\Lock} - P^{mkt}_{t_e})$  & first-order proxy for the T-Lock payoff \\
\eqref{eq:fwd-repo}                           & $\Fwdpx{t}{t_e}$                                   & forward bond price under repo (zero haircut) \\
\eqref{eq:delta},\eqref{eq:gamma-local}       & $\Delta$, $\Gam$                                   & T-Lock first and second derivatives \\
\bottomrule
\end{tabular}
\smallskip

\textbf{Invariants and identities to verify:}
\begin{enumerate}[noitemsep]
  \item $\Bond{t}{\IRR_t} = P^{mkt}_t$ holds at the canonical case
        ($P^{mkt} = 104.1055$, $\IRR = 2.717\%$) within float tolerance.
  \item Eq.~\eqref{eq:fwd-repo} evaluated at the canonical inputs equals
        $104.74$ (dirty) within $\pm 0.01$.
  \item \emph{Overhedge inequality}: for any $\Lock$, any $y$ in a
        reasonable yield range, and any positive-coupon schedule,
        \[
          \Bond{}{\Lock} + \Diff{y}{1}[\Bond{}{\cdot}](y - \Lock) \;\le\; \Bond{}{y}.
        \]
        A global property, following from $\Diff{y}{2}[\Bond{}{\cdot}] > 0$.
  \item \emph{Gamma sign (local)}: for $a=1$ and $y$ in a neighbourhood
        of $\Lock$, $\Diff{y}{2}[g] < 0$. Sharp threshold given by
        \eqref{eq:gamma-cond}.
  \item \emph{Closed-form vs.\ numerical agreement}: the closed-form
        derivative \eqref{eq:bond-derivs} equals a numerical derivative
        of $\Bond{t}{y}$ at any test yield.
  \item \emph{Monotonicity in repo}: $\Fwdpx{t}{t_e}$ in
        \eqref{eq:fwd-repo} is strictly increasing in $\Repo$ (under
        the standing assumption that the spot price exceeds the present
        value of intermediate coupons --- true here).
\end{enumerate}
\end{validation}

% =====================================================================
\section{Glossary of symbols (paper's notation)}

\begin{tabular}{@{}p{3cm}p{12cm}@{}}
\toprule
$P_t(y)$        & Bond dirty price at $t$ as a function of yield $y$ \\
$P_t^{mkt}$     & Observed market dirty price ($= P_t(\IRR_t)$) \\
$\IRR_t$        & Internal rate of return at $t$ \\
$L$             & Locked yield (T-Lock strike) \\
$\Riskf_t$      & $-\partial P_t/\partial y$ at $y=\IRR_t$; equivalently $-10^4 \cdot \DVone$ \\
$\DVone$        & Price change per 1~bp positive yield shift \\
$g$             & T-Lock payoff at expiry \\
$a$             & $+1$ (long) or $-1$ (short) unitary T-Lock \\
$t_e$           & T-Lock expiry \\
$D_{tt_e}$      & Risk-free discount factor from $t$ to $t_e$ \\
$D_y^k[\,\cdot\,]$ & Euler shorthand for $\partial^k/\partial y^k$ \\
$P_t^{(0)}$     & Bond price stripped of coupons (coupon $c=0$); pure-discount component \\
$A_t(y)$        & Coupon annuity $\sum_{t_i>t} \alpha_i e^{-(t_i-t)y}$ (CC form) \\
$\alpha_i$      & Year fraction for accrual period ending at $t_i$ \\
$r^{\mathrm{repo}}_{uv}$ & Simply-compounded forward repo rate, $u \to v$ \\
$r^{\mathrm{fun}}$ & Unsecured funding rate \\
$h^{\mathrm{cut}}$ & Repo haircut, $\in [0,1]$ \\
$\mathrm{ForwardPrice}_t(t_e)$ & $t_e$-forward bond price at $t$ \\
$R_t$ ($\widehat{R}_t$) & IRR of the current (new) on-the-run benchmark at $t$ \\
$t_{\mathrm{roll}}$ & Date of benchmark roll in a Then-Current-T-Lock \\
\bottomrule
\end{tabular}

% =====================================================================
\begin{todobox}
\begin{itemize}[noitemsep]
  \item Reproduce Fig.~1 (T-Lock payoff vs.\ forward overhedge) on the
        24-Jan-2019 test case. Provides a visual smoke test for the
        dirty-price and forward-price formulas.
  \item Reproduce Fig.~2 \& 3 (Delta, Gamma plots) --- shows the
        sign-flip at extreme yields predicted by \eqref{eq:delta-cond}
        and \eqref{eq:gamma-cond}.
  \item Section 4 statistical analysis: implement the roll-P\&L
        estimator and check against paper's Table 2 (IRR statistics
        from 2008 Treasury auction history). Lower priority since
        this is empirical, not formula-validation.
  \item Repo-special scenario: re-price the test case with $\Repo$ shifted
        $-300$~bps (paper p.~12) to confirm forward price moves in the
        documented direction.
  \item Open question: paper assumes a flat short-term repo curve in
        Sec.~2.4 (around eq.~50). How sensitive is the analysis to
        the curve shape? Should the validation suite include a
        non-flat repo curve case?
\end{itemize}
\end{todobox}

\end{document}
